\chapter{Extensions of the Inventory Framework}\label{ch:hf:extensions-of-the-inventory-framework}

The exact solution of chapter 3 tells a market maker holding four units to offer them below the mid. Real market makers stop buying instead:
they bound the inventory. They also hold several correlated books at once, they know something about where the price is going, and they know
that the next trade is more likely to come from someone who knows more than they do. Each of these enters the same linear system. In this
chapter's simulations, quoting two assets whose prices are correlated 0.8 as one book instead of two carries 21\% less inventory risk for the same
expected profit; a market maker who knows the drift holds the position the formula targets; and one whose fills are followed by an adverse move
does best by widening each quote by half that move, not by all of it.

\section{Inventory bounds and the asymptotic solution}

\begin{definition}[Inventory bound]\label{def:hf:extensions-of-the-inventory-framework:bound}
An \emph{inventory bound}\index{inventory bound} $\bar q$ is a hard limit on the market maker's position: at $q=\bar q$ it does not bid, at
$q=-\bar q$ it does not offer, so that $|q_t|\le\bar q$ at all times.
\end{definition}

The bound is part of the state space of \cref{prop:hf:inventory-models:cj} of chapter 3, which is why that solution was exact on a finite grid. Its
first use is practical: whatever the model says, the position cannot run away. Its second is analytical: on a finite grid the long-horizon
behaviour of the solution is an eigenvector problem.

\begin{proposition}[The stationary quotes]\label{prop:hf:extensions-of-the-inventory-framework:stationary}
As $T-t\to\infty$, the depths of chapter 3's running-penalty model converge to time-independent depths given by the same formulas with
$h=\tfrac1k\ln v$, where $v$ is the eigenvector of the largest eigenvalue of $M$.
\end{proposition}

\begin{proof}
$M$ is symmetric with positive off-diagonal entries next to the diagonal, so its largest eigenvalue $\mu_1$ is simple and its eigenvector $v$ can be
taken positive (Perron--Frobenius). Then $\omega(t)=e^{M(T-t)}z=\sum_ie^{\mu_i(T-t)}(v_i^\top z)v_i$ is dominated by the first term as $T-t\to\infty$,
since $v^\top z>0$. The depths depend on $\omega$ only through ratios across $q$, so the factor $e^{\mu_1(T-t)}(v^\top z)$ drops out.
\end{proof}

In the chapter 3 parameters ($A=140$, $k=1.5$, $\phi=0.5$) a horizon of one unit of time is already long: the depths at its start agree with the
stationary ones to eight decimals. Guéant, Lehalle and Fernandez-Tapia solved the exponential-utility version the same way and approximated the
eigenvector in closed form. The same approximation for the running penalty gives a formula simple enough to put in a pricing function.

\begin{proposition}[Closed-form stationary quotes]\label{prop:hf:extensions-of-the-inventory-framework:closed}
If the eigenvector is approximately Gaussian, $v_q\propto e^{-\beta q^2/2}$ with $\beta$ small, then $\beta^2=\phi ke/A$ and, with
$c=\beta/k=\sqrt{\phi e/(Ak)}$,
\[
\delta^b(q)\simeq\frac1k+\frac{2q+1}{2}\,c,\qquad \delta^a(q)\simeq\frac1k-\frac{2q-1}{2}\,c .
\]
\end{proposition}

\begin{proof}
The eigenvalue equation is $Ae^{-1}(v_{q-1}+v_{q+1})-\phi kq^2v_q=\mu_1v_q$. With the Gaussian ansatz, $v_{q\pm1}/v_q=e^{-\beta(\pm2q+1)/2}$, whose sum
is $2e^{-\beta/2}\cosh(\beta q)\simeq2-\beta+\beta^2q^2$ to second order. Matching the terms in $q^2$ gives $Ae^{-1}\beta^2=\phi k$. Then
$h(q)-h(q\pm1)=\tfrac1k\ln(v_q/v_{q\pm1})=\tfrac{\beta}{2k}(\pm2q+1)$.
\end{proof}

Each unit of inventory shifts both quotes by $c$, and the spread stays $2/k+c$. For $\gamma\to0$ with $\phi=\tfrac12\gamma\sigma^2$, the
Guéant--Lehalle--Fernandez-Tapia formula reduces to this one. The approximation needs $\beta$ small, that is a penalty small against the fill rate:

\begin{center}
\small
\begin{tabular}{@{}lccccc@{}}
\toprule
bid depth, $\phi=0.5$ & $q=0$ & $q=1$ & $q=2$ & $q=3$ & worst gap, $|q|\le10$\\
\midrule
exact, stationary & 0.708 & 0.790 & 0.872 & 0.953 & \\
closed form & 0.707 & 0.787 & 0.868 & 0.948 & 0.034\\
\midrule
worst gap for $\phi=5$ / $\phi=50$ & & & & & 0.70 / 5.0\\
\bottomrule
\end{tabular}
\end{center}

\Cref{fig:hf:extensions-of-the-inventory-framework:closed} shows the failure at $\phi=5$: the exact depths curve at large inventories, where the
eigenvector is far from Gaussian; the straight line does not.

\begin{omfigure}
\begin{tikzpicture}
\begin{axis}[width=11cm,height=5.2cm,xlabel={\small inventory $q$},ylabel={\small bid depth},tick label style={font=\footnotesize},
  xmin=-8,xmax=8,grid=major,grid style={gray!25},table/col sep=comma,legend style={font=\scriptsize,at={(0.5,-0.3)},anchor=north},legend columns=4]
\addplot[omDef,thick] table[x=q,y=exact05]{figdata/hft/04-extensions-of-the-inventory-framework/closed.csv};
\addplot[omDef,thick,dashed] table[x=q,y=closed05]{figdata/hft/04-extensions-of-the-inventory-framework/closed.csv};
\addplot[omThm,thick] table[x=q,y=exact5]{figdata/hft/04-extensions-of-the-inventory-framework/closed.csv};
\addplot[omThm,thick,dashed] table[x=q,y=closed5]{figdata/hft/04-extensions-of-the-inventory-framework/closed.csv};
\legend{{exact, $\phi=0.5$},{closed form, $\phi=0.5$},{exact, $\phi=5$},{closed form, $\phi=5$}}
\end{axis}
\end{tikzpicture}

\omcaption{Stationary bid depth against inventory: exact (eigenvector) and closed form, for two running penalties ($A=140$, $k=1.5$, inventory
within $\pm20$). Data: \texttt{firm.multimm}.}\label{fig:hf:extensions-of-the-inventory-framework:closed}
\end{omfigure}

\section{Many assets and correlated inventory}

\begin{definition}[Multi-asset market making]\label{def:hf:extensions-of-the-inventory-framework:multi}
\emph{Multi-asset market making}\index{multi-asset market making} quotes several instruments from one inventory vector $q\in\mathbb Z^d$ and one
risk charge on the whole book, $\phi\,q^\top\Sigma q$ per unit of time, with $\Sigma$ the covariance of the instruments' price changes, so that each
quote depends on the inventory of every instrument.
\end{definition}

\begin{proposition}[Exact multi-asset quotes]\label{prop:hf:extensions-of-the-inventory-framework:multi}
If the instruments share $A$ and $k$, the substitution $h=\tfrac1k\ln\omega$ turns the problem on the grid $\{-\bar q,\dots,\bar q\}^d$ into the
linear system $\partial_t\omega+M\omega=0$, with $M_{q,q}=-\phi k\,q^\top\Sigma q$ and $Ae^{-1}$ between neighbouring inventories $q\pm e_i$; the
depths of instrument $i$ are $\delta^{b,a}_i(q)=\tfrac1k+h(q)-h(q\pm e_i)$, and \cref{prop:hf:extensions-of-the-inventory-framework:stationary}
holds on the grid.
\end{proposition}

\begin{proof}
As in chapter 3, instrument by instrument: each side's supremum contributes
\[
\tfrac1k\,Ae^{-1}\,e^{k(h(q\pm e_i)-h(q))},
\]
which is linear in $\omega$ after the substitution because every instrument has the same $k$.
\end{proof}

The shared $k$ is the price of exactness; with different $k_i$ the system is no longer linear and the multi-asset literature (Guéant 2017) turns to
quadratic approximations of the value function. The grid has $(2\bar q+1)^d$ points, 289 for two instruments with $\bar q=8$: fine for a pair, not
for a hundred names.

Correlation creates a cross-skew. Under the joint policy (\cref{fig:hf:extensions-of-the-inventory-framework:cross}), a market maker short four
units of the second instrument bids for the first at 0.572 from the mid when flat in it, against 0.700 when flat in both: a long position in the
first hedges the short in the second. Asset by asset, the bid would be 0.700 whatever the other position.

\begin{omfigure}
\begin{tikzpicture}
\begin{axis}[width=11cm,height=5.2cm,xlabel={\small inventory of asset 1},ylabel={\small bid depth of asset 1},tick label style={font=\footnotesize},
  xmin=-6,xmax=6,grid=major,grid style={gray!25},table/col sep=comma,legend style={font=\scriptsize,at={(0.5,-0.3)},anchor=north},legend columns=3]
\addplot[omDef,thick,mark=*,mark size=1pt] table[x=q1,y=short4]{figdata/hft/04-extensions-of-the-inventory-framework/cross.csv};
\addplot[black,thick,mark=square*,mark size=1pt] table[x=q1,y=flat]{figdata/hft/04-extensions-of-the-inventory-framework/cross.csv};
\addplot[omThm,thick,mark=triangle*,mark size=1.3pt] table[x=q1,y=long4]{figdata/hft/04-extensions-of-the-inventory-framework/cross.csv};
\legend{{short 4 of asset 2},{flat in asset 2},{long 4 of asset 2}}
\end{axis}
\end{tikzpicture}

\omcaption{The joint policy's bid depth for asset 1 against its own inventory, for three inventories of asset 2 (two assets with price
correlation 0.8, penalty $\phi=0.1$ on $q^\top\Sigma q$, $A=140$, $k=1.5$, $\sigma=2$). Data: \texttt{hf\_extensions.cross\_skew}.}\label{fig:hf:extensions-of-the-inventory-framework:cross}
\end{omfigure}

The cross-skew pays. On 1\,000 simulated paths of five units of time, with the two mids correlated 0.8, the joint policy and the asset-by-asset
policy (the same penalty applied to each asset's own variance) earn almost the same, but the joint book carries less inventory risk, measured as the
time average of $q^\top\Sigma q$ (\cref{fig:hf:extensions-of-the-inventory-framework:frontier}). At the joint policy's mean profit for $\phi=0.1$,
673.3, the asset-by-asset frontier carries 30.8 units of inventory risk and the joint policy 24.3: 21\% less. Across the six penalties the
reduction runs from 10\% to 22\%.

\begin{omfigure}
\begin{tikzpicture}
\begin{axis}[width=11cm,height=5.2cm,xlabel={\small inventory risk (time average of $q^\top\Sigma q$)},ylabel={\small mean profit},
  tick label style={font=\footnotesize},grid=major,grid style={gray!25},table/col sep=comma,
  legend style={font=\scriptsize,at={(0.5,-0.3)},anchor=north},legend columns=2]
\addplot[omThm,thick,mark=*,mark size=1.3pt] table[x=risk_joint,y=mean_joint]{figdata/hft/04-extensions-of-the-inventory-framework/frontier2.csv};
\addplot[omDef,thick,dashed,mark=square*,mark size=1.3pt] table[x=risk_sep,y=mean_sep]{figdata/hft/04-extensions-of-the-inventory-framework/frontier2.csv};
\legend{joint book,asset by asset}
\end{axis}
\end{tikzpicture}

\omcaption{Mean profit against inventory risk for two assets correlated 0.8, quoted as one book or asset by asset, for penalties $\phi$ from 0.01
(right) to 0.5 (left); 1\,000 common paths of five units of time. Data: \texttt{hf\_extensions.frontier2}.}\label{fig:hf:extensions-of-the-inventory-framework:frontier}
\end{omfigure}

\section{Signals and drift inside the control problem}

\begin{definition}[Drift-adjusted quoting]\label{def:hf:extensions-of-the-inventory-framework:drift}
\emph{Drift-adjusted quoting}\index{drift-adjusted quoting} solves the market maker's control problem with the forecast drift $\mu$ of the mid
inside it, so that the forecast shifts the quotes through the value of holding inventory rather than being added to them as a separate skew.
\end{definition}

With $dS_t=\mu\,dt+\sigma\,dW_t$, holding $q$ earns $\mu q$ per unit of time, and the running reward is $\mu q-\phi q^2=-\phi(q-q_0)^2+\text{const}$ with
$q_0=\mu/(2\phi)$.

\begin{proposition}[A drift sets a target position]\label{prop:hf:extensions-of-the-inventory-framework:target}
With a constant drift $\mu$, the diagonal of $M$ becomes $k(\mu q-\phi q^2)$, and in the closed form of
\cref{prop:hf:extensions-of-the-inventory-framework:closed} the market maker quotes as if its inventory were $q-q_0$: it targets the position
$q_0=\mu/(2\phi)$.
\end{proposition}

\begin{proof}
The running reward enters $M$'s diagonal multiplied by $k$, as $-\phi q^2$ did. Completing the square, the eigenvalue equation of
\cref{prop:hf:extensions-of-the-inventory-framework:closed} holds with $q$ replaced by $q-q_0$ up to a constant absorbed in $\mu_1$.
\end{proof}

Guéant, Lehalle and Fernandez-Tapia's formula with a drift has the same shape, a shift of the inventory by $\mu/(\gamma\sigma^2)$, which is $q_0$ for
$\phi=\tfrac12\gamma\sigma^2$. Cartea and Wang put a stochastic alpha signal inside the problem in the same spirit. In the simulation with $\mu=2$,
$\phi=0.5$ (so $q_0=2$), the drift-adjusted market maker holds 1.97 units on average and earns 355.6, against 335.7 for the market maker that ignores the
drift and stays flat: 5.9\% more, for a standard deviation of 18.0 instead of 15.8. The forecast is worth money only through the position it makes
the market maker hold; chapter 7 builds such forecasts at the horizon of seconds.

\section{Adverse selection inside the model}

Adverse selection enters the model in two ways, and they call for different widenings.

\begin{proposition}[Widen by the cost, or by half the move]\label{prop:hf:extensions-of-the-inventory-framework:adverse}
(i) If each fill costs an expected $\varepsilon$ (the fill price is on average $\varepsilon$ worse than the mid), the off-diagonal of $M$ becomes
$Ae^{-1-k\varepsilon}$ and both optimal depths increase by $\varepsilon$ at every inventory. (ii) If instead the mid moves permanently by $\xi$ against
the market maker after each of its fills, the optimal depths are those of (i) with $\varepsilon=\xi/2$.
\end{proposition}

\begin{proof}
(i) Each supremum becomes $\sup_\delta Ae^{-k\delta}(\delta-\varepsilon+\Delta h)$: the optimum is $1/k+\varepsilon-\Delta h$ and its value carries a
factor $e^{-k\varepsilon}$. (ii) After a bid fill the whole position $q+1$ is revalued by $-\xi$, after an ask fill $q-1$ by $+\xi$: the brackets are
$\delta-\xi(q+1)+h(q+1)-h(q)$ and $\delta+\xi(q-1)+h(q-1)-h(q)$. With $h(q)=\tilde h(q)+\tfrac12\xi q^2$ both become $\delta-\tfrac12\xi+\Delta\tilde
h$, the problem (i) for $\tilde h$ with $\varepsilon=\xi/2$, and the depths expressed through $\tilde h$ are those of (i).
\end{proof}

Half, because the move that hurts the unit just traded also revalues the rest of the position, and on average the position the fill leaves is on the
right side of the move as often as on the wrong one. The simulation of (ii) with $\xi=0.3$, $\phi=0.5$ agrees, provided the time step is fine
enough:

\begin{center}
\small
\begin{tabular}{@{}lccc@{}}
\toprule
widening of each depth & none & $\xi/2=0.15$ & $\xi=0.3$\\
\midrule
mean profit & 262.7 & 267.2 & 258.9\\
profit less the penalty & 253.1 & 258.6 & 251.2\\
fills & 498 & 396 & 315\\
profit less the penalty, time step four times coarser & 270.8 & 270.2 & 259.3\\
\bottomrule
\end{tabular}
\end{center}

With the coarse step (0.005, when a flat quote has about a 24\% chance of a fill per step), the unadjusted policy looked marginally best; divided by four, the step
reproduces the proposition. The rule of One Quant Book 7 applies: a simulation that a named result leans on is rerun with a quarter of the time
step.

\section{Tutorial: one linear system, four extensions}\label{tut:hf:extensions-of-the-inventory-framework:tut}

\begin{tutorial}
\textbf{Goal.} Compute stationary, closed-form, multi-asset, drift-adjusted and adverse-selection-adjusted quotes from one construction, and test
each in simulation. \textbf{End state:} the three tables and three figures of the chapter.
\begin{tutsteps}
\item \textbf{Stationary and closed form.} The Perron eigenvector of $M$ gives the long-horizon depths; the Gaussian approximation gives the
formula (\cref{lst:hf:extensions-of-the-inventory-framework:stationary}).
\omcode{code/firm/multimm/firm_multimm.py}{49}{69}{Stationary depths from the eigenvector, and their closed form.}\label{lst:hf:extensions-of-the-inventory-framework:stationary}
\item \textbf{Two assets.} The same construction on the grid of inventory pairs
(\cref{lst:hf:extensions-of-the-inventory-framework:multi}); \texttt{simulate\_multi} draws correlated mids.
\omcode{code/firm/multimm/firm_multimm.py}{76}{94}{The multi-asset linear system and its eigenvector.}\label{lst:hf:extensions-of-the-inventory-framework:multi}
\item \textbf{Frontiers.} \texttt{hf\_extensions.frontier2()} runs both policies for six penalties on 1\,000 common paths;
\texttt{equal\_capture()} interpolates the asset-by-asset frontier at each joint mean.
\item \textbf{Drift and adverse selection.} \texttt{drift\_case(adjust)} and \texttt{adverse\_case(adjust, dt)}.
\end{tutsteps}
\textbf{What to change next.} Give the two assets different fill rates and compare the grid solution with a quadratic approximation; make the drift
switch sign at random (a Markov signal) and add it to the state.
\end{tutorial}

\section{Build: the multi-asset and signal extensions}\label{bld:hf:extensions-of-the-inventory-framework:multimm}

\begin{build}
\bfield{Purpose} Long-horizon quotes that a production quoter can evaluate in microseconds (a table or a formula), for one or several instruments,
with a drift and adverse selection inside the problem.
\bfield{Interface} \texttt{firm.multimm}: \texttt{stationary(A, k, phi, qmax, mu, eps)}, \texttt{asymptotic(A, k, phi, q, mu, eps)},
\texttt{MultiStationary(A, k, phi, cov, qmax).depths(q)}, \texttt{simulate\_multi(policy, T, dt, cov, A, k, paths, seed, qmax)},
\texttt{table\_policy}; \texttt{firm.invmm.CJSolution(\dots, mu, eps)} for finite horizons.
\bfield{Rules} The \omterm{def:hf:extensions-of-the-inventory-framework:bound}{inventory bound} is never breached. A permanent post-fill move $\xi$ is passed as $\varepsilon=\xi/2$. Simulations that support a
printed result are repeated with a quarter of the time step.
\bfield{Acceptance tests} \texttt{code/firm/multimm/tests/}: the stationary depths equal the finite-horizon ones at the start of a long horizon; the
closed form within 0.05 of the exact depths for a small penalty and its per-unit skew equal to $c$; a drift shifts the quotes by $q_0$; $\varepsilon$
shifts both closed-form depths by $\varepsilon$; two uncorrelated assets quote like two single assets; with correlation, a short position in one
asset lowers the bid depth of the other; the simulator is deterministic and bounded.
\bfield{Stretch} Different $k_i$ with a quadratic approximation of $h$; a Markov-modulated drift in the state.
\end{build}

\begin{omsources}
\item O.~Guéant, C.-A.~Lehalle and J.~Fernandez-Tapia, ``Dealing with the inventory risk: a solution to the market making problem'',
\emph{Mathematics and Financial Economics} 7(4), 2013.
\item O.~Guéant, ``Optimal market making'', \emph{Applied Mathematical Finance} 24(2), 2017.
\item \'A.~Cartea and Y.~Wang, ``Market making with alpha signals'', \emph{International Journal of Theoretical and Applied Finance} 23(3), 2020.
\end{omsources}

\section{Exercises}

\begin{exercise}[$\star$]\label{exo:hf:extensions-of-the-inventory-framework:1}
Compute $c=\sqrt{\phi e/(Ak)}$ for $A=140$, $k=1.5$, $\phi=0.5$, and the closed-form bid and ask depths at $q=2$.
\end{exercise}

\begin{exercise}[$\star$]\label{exo:hf:extensions-of-the-inventory-framework:2}
A market maker with $\phi=0.5$ forecasts a drift of $\mu=1$. What position does it target?
\end{exercise}

\begin{exercise}[$\star$]\label{exo:hf:extensions-of-the-inventory-framework:3}
How many grid points does the exact multi-asset solution need for three instruments with $\bar q=8$, and for ten?
\end{exercise}

\begin{exercise}[$\star\star$]\label{exo:hf:extensions-of-the-inventory-framework:4}
Show that the closed-form spread $\delta^a+\delta^b$ does not depend on inventory, and give it for the parameters of exercise 1.
\end{exercise}

\begin{exercise}[$\star\star$]\label{exo:hf:extensions-of-the-inventory-framework:5}
Why does the joint policy bid more aggressively for asset 1 when short asset 2, and would it do so if the correlation were negative?
\end{exercise}

\begin{exercise}[$\star\star$]\label{exo:hf:extensions-of-the-inventory-framework:6}
Explain in words why a permanent post-fill move $\xi$ calls for widening by $\xi/2$ rather than $\xi$.
\end{exercise}

\begin{exercise}[$\star\star\star$]\label{exo:hf:extensions-of-the-inventory-framework:7}
\emph{Coding.} With \texttt{stationary(140, 1.5, 0.5, 20, mu=2.0)}, give the bid and ask depths at $q=0$ and at $q=2$, and compare them with the
depths at $q=0$ without a drift.
\end{exercise}

\begin{exercise}[$\star\star\star$]\label{exo:hf:extensions-of-the-inventory-framework:8}
\emph{Find the flaw.} ``Our simulation shows that widening for adverse selection loses money, so we quote at the monopoly depth and ignore it.''
\end{exercise}

\section{Problem: Two Books, One Risk}

\begin{problem}[{Weekend problem --- two books, one risk}]\label{pb:hf:extensions-of-the-inventory-framework:1}
A desk makes markets in two instruments that move together and has a signal for one of them.

\medskip\noindent\textbf{Part I --- Bounds and the long horizon.}
\begin{enumerate}
\item Define an \omterm{def:hf:extensions-of-the-inventory-framework:bound}{inventory bound} and say why it makes the long-horizon quotes an
eigenvector problem.
\item Prove \cref{prop:hf:extensions-of-the-inventory-framework:stationary}.
\item Derive the closed form of \cref{prop:hf:extensions-of-the-inventory-framework:closed}.
\item When is the closed form accurate? Give the worst gaps for $\phi=0.5$, 5 and 50.
\end{enumerate}

\medskip\noindent\textbf{Part II --- Two assets.}
\begin{enumerate}[resume]
\item Define \omterm{def:hf:extensions-of-the-inventory-framework:multi}{multi-asset market making} and write the risk charge.
\item Why must the instruments share $k$ for the exact solution?
\item Give the cross-skew: asset 1's bid depth at $q_1=0$ when short four, flat and long four of asset 2.
\item Describe the simulation and its measure of inventory risk.
\item State the \emph{named result}: the reduction in inventory risk from quoting the two assets as one book, at equal mean profit, for $\phi=0.1$,
and its range across penalties.
\item Why is the reduction in P\&L standard deviation much smaller than in inventory risk?
\end{enumerate}

\medskip\noindent\textbf{Part III --- The signal.}
\begin{enumerate}[resume]
\item Define \omterm{def:hf:extensions-of-the-inventory-framework:drift}{drift-adjusted quoting}.
\item Prove \cref{prop:hf:extensions-of-the-inventory-framework:target}.
\item Give the simulated gain and average position with $\mu=2$, $\phi=0.5$.
\item What is the drift worth if the market maker's bound is below $q_0$?
\end{enumerate}

\medskip\noindent\textbf{Part IV --- Adverse selection.}
\begin{enumerate}[resume]
\item Prove both parts of \cref{prop:hf:extensions-of-the-inventory-framework:adverse}.
\item Give the simulated profit and penalised objective for no widening, $\xi/2$ and $\xi$.
\item Why did the coarse time step point to no widening?
\item Which of the chapter's four extensions would you implement first in a production quoter, and why?
\item How would you estimate $\xi$ from your own fills?
\item In one sentence: what do bounds, correlation, signals and adverse selection have in common in this framework?
\end{enumerate}
\end{problem}

\section{Interview questions}

\begin{interviewq}[$\star$ \iqroles{trader}]\label{iq:hf:extensions-of-the-inventory-framework:1}
You make markets in an index future and its exchange-traded fund. You are long the fund. How should that change your quotes in the future?
\end{interviewq}

\begin{interviewq}[$\star\star$ \iqroles{researcher}]\label{iq:hf:extensions-of-the-inventory-framework:2}
Why do the finite-horizon quotes stop depending on time far from the horizon?
\end{interviewq}

\begin{interviewq}[$\star\star$ \iqroles{researcher}]\label{iq:hf:extensions-of-the-inventory-framework:3}
Your fills are followed by a 0.3-tick adverse move on average. By how much do you widen, and what do you need to know first?
\end{interviewq}

\begin{interviewq}[$\star\star$ \iqroles{developer}]\label{iq:hf:extensions-of-the-inventory-framework:4}
A quoter must evaluate its depths in under a microsecond. How do you deliver a model whose exact solution is an eigenvector problem?
\end{interviewq}

\begin{interviewq}[$\star\star$ \iqroles{risk}]\label{iq:hf:extensions-of-the-inventory-framework:5}
A desk measures inventory risk instrument by instrument. What does it miss, and how would you measure it instead?
\end{interviewq}

\begin{interviewq}[$\star\star\star$ \iqroles{researcher}]\label{iq:hf:extensions-of-the-inventory-framework:6}
In a running-penalty market-making model, show that a constant drift $\mu$ makes the market maker target the position $\mu/(2\phi)$.
\end{interviewq}
